% =====================================================================
%  第 20 章  重夸克物理：B 介子衰变与重夸克偶素            （ch20_heavy.tex，writer W23）
%  约定：本书约定（WRITER_SPEC_v2 §3）；n_A=(1^+,0^-,0_T)、n_B=(0^+,1^-,0_T)（不用 SCET 的 n、nbar）；
%  D_mu = d_mu - i g A_mu^a t^a；D = 4-2 eps，S_eps=(4pi)^eps e^{-eps gamma_E}；
%  BBL 与 Manohar--Wise 用 D = d + i g A（g -> -g）。验证脚本：writing/work/W23/s1_bdpi.py … s8_review.py（日志在 logs/）。
% =====================================================================
\chapter{重夸克物理：$B$ 介子衰变与重夸克偶素}\label{ch:heavy}

重夸克质量 $m_Q\gg\Lqcd$ 本身就是硬标度：$b$ 夸克衰变使末态轻介子以近光速运动，出现与硬遍举过程（\cref{ch:excl}）相同的共线结构；重夸克偶素中的相对速度 $v\ll1$ 提供第二个展开参数。本章只有一条严格结果——电磁衰变匹配系数的树图值与单圈值（[R, $R_{\mathrm{PT}}$]）；其余结果的地位从全阶论证到猜想不等，见下框。贯穿全章的一点是：衰变率是固定物理点上的 Minkowski 量，把这些结果用于固定运动学点上的强子衰变率还需要局域夸克—强子对偶（\cref{conj:ee-duality}(b)，[C]），它不在任何微扰证明的范围之内。

\begin{tcolorbox}[enhanced,breakable,title=本章证明了什么／没有证明什么,fonttitle=\bfseries,
  colback=gray!3!white,colframe=gray!55!black,boxrule=0.6pt,arc=2pt,
  left=5pt,right=5pt,top=3pt,bottom=3pt,before skip=6pt,after skip=6pt]
\textbf{证明了：}\cref{prop:heavy-nrqcd-matching}（\tagR，$R_{\mathrm{PT}}$：电磁衰变匹配系数的树图值与对 $\varepsilon$ 精确的单圈 $c_v$）。\par
\textbf{给出全阶论证：}\cref{thm:heavy-bdpi,prop:heavy-btopi-ff}（\tagA；瓶颈：\cref{rem:scet-assumptions}的 SCET 前提，$B\to D\pi$ 另需端点压低假设）；\cref{thm:heavy-inclusive-b}（\tagA；瓶颈：微扰 OPE \cref{thm:dis-ope}，端点区另需 SCET 前提）；\cref{thm:heavy-nrqcd-decay}（\tagA；瓶颈：拓扑因子化的全阶形式）；\cref{prop:heavy-quarkonium-fragmentation}（\tagA；瓶颈：强子对撞中有色末态的 Glauber 相消）。强子层次均另需\cref{asm:intro-bridge}。\par
\textbf{没有证明（或已知不成立）：}两个轻介子末态领头幂次因子化的全阶形式（\cref{prop:heavy-qcdf-lp}，\tagP）；其幂次修正的标准共线因子化——显式计算给出发散的卷积（\cref{prop:heavy-qcdf-endpoint}，\tagV）；产生的 NRQCD 因子化（\cref{conj:heavy-nrqcd-production}，\tagC）；固定运动学点上所需的局域对偶（\cref{conj:ee-duality}，\tagC）。
\end{tcolorbox}

\paragraph{记号。}
本章用 $n_A=(1^+,0^-,\mathbf 0_T)$、$n_B=(0^+,1^-,\mathbf 0_T)$，引用 SCET 文献（\cref{ch:scet}）时作换算 $n\to\sqrt2\,n_A$、$\bar n\to\sqrt2\,n_B$。被发射的轻介子或喷注总取为沿 $+z$ 运动，$p\simeq p^+n_A$：其组分经 $n_A\cdot A_s=A_s^-$ 与软胶子耦合，与之共线的胶子求和成沿 $n_B$ 的 Wilson 线。$\lambda=\Lqcd/Q$，$Q$ 为硬标度。重夸克有效理论（HQET）的场为 $h_v=\ee^{\ii m_Qv\cdot x}\frac{1+\sla v}{2}Q$，$v$ 为强子的四速度。BBL\cite{BodwinBraatenLepage1995}与 Manohar--Wise\cite{ManoharWise1994}用 $D_\mu=\partial_\mu+\ii gA_\mu$，其公式已按 $g\to-g$ 换算到\cref{def:prelim-qcd}的约定；含场强的式子一律用协变导数的对易子写出。

% ---------------------------------------------------------------------
\section{$B\to D\pi$：领头幂次的全阶因子化}\label{sec:heavy-bdpi}

在“I 类”衰变 $\bar B^0\to D^+\pi^-$ 中（\cref{fig:heavy-bdpi}），弱顶点 $b\to c\,d\bar u$ 直接产生 $\pi^-$，旁观反夸克 $\bar d$ 进入 $D^+$。按色透明性（colour transparency）\cite{Bjorken:1988kk}，快速的 $\pi$ 形成之初是尺度 $\sim1/E_\pi$ 的色单态小偶极子，软胶子分辨不出其中的色荷，振幅近似朴素因子化\cite{Dugan:1990de,Politzer:1991au}；Beneke、Buchalla、Neubert、Sachrajda（BBNS）在两圈显式验证了这一点\cite{BenekeEtAl2000}，Bauer、Pirjol、Stewart（BPS）在软—共线有效理论（SCET）与 HQET 中给出了对 $\as$ 所有阶的论证\cite{BauerPirjolStewart2001}。

%%% BEGIN fig:heavy-bdpi
\begin{figure}[htbp]
\centering
\begin{tikzpicture}[fd,
  fd2 heavy/.style={draw,solid,line width=1.3pt,fd arrow=#1},
  fd2 heavy/.default=0.5,
  fd2 meson/.style={bound state,line cap=butt},
  fd2 weak/.style={rectangle,fill=black,draw=none,inner sep=0pt,minimum size=5.6pt},
  fd2 soft/.style={gluon,draw=green!45!black},
  fd2 cross/.pic={\draw[red!75!black,line width=1.1pt,line cap=round] (-0.12,-0.12)--(0.12,0.12) (-0.12,0.12)--(0.12,-0.12);},
  fd2 bracket/.style={draw=gray!55!black,line width=0.5pt},
  fd2 lab/.style={font=\small},
  fd2 note/.style={font=\footnotesize}]
  \coordinate (W) at (0,0.45);
  \coordinate (dq) at (1.449,2.820);   % end of the d line (inside the pion blob)
  \coordinate (uq) at (1.851,2.558);   % end of the u-bar line (inside the pion blob)
  % ---- mesons (double lines) ----
  \draw[fd2 meson] (-5.15,0) -- (-3.9,0);
  \draw[fd2 meson] (3.9,0) -- (5.15,0);
  \draw[fd2 meson] (2.36,3.78) -- (2.84,4.51);
  % ---- soft gluons from the B-D system to the pion pair (decouple) ----
  \coordinate (sd) at ($(W)!0.42!(dq)$);
  \coordinate (su) at ($(W)!0.36!(uq)$);
  \draw[fd2 soft] (-3.42,0.72) -- (sd);
  \draw[fd2 soft] (3.42,0.72) -- (su);
  \pic at ($(-3.42,0.72)!0.5!(sd)$) {fd2 cross};
  \pic at ($(3.42,0.72)!0.5!(su)$) {fd2 cross};
  % ---- quark lines (arrows = fermion-number flow) ----
  \draw[fd2 heavy=0.55] (-3.4,0.45) -- (W);
  \draw[fd2 heavy=0.55] (W) -- (3.4,0.45);
  \draw[quark] (3.4,-0.45) -- (-3.4,-0.45);
  \draw[quark=0.72] (W) -- (dq);
  \draw[quark=0.6] (uq) -- (W);
  % ---- blobs on top ----
  \node[hadron blob,ellipse,minimum width=0.8cm,minimum height=1.8cm,inner sep=0pt] at (-3.6,0) {};
  \node[hadron blob,ellipse,minimum width=0.8cm,minimum height=1.8cm,inner sep=0pt] at (3.6,0) {};
  \node[jet blob,ellipse,minimum width=1.6cm,minimum height=0.72cm,inner sep=0pt,rotate=57] at (1.95,3.15) {};
  \node[fd2 weak] at (W) {};
  % ---- labels ----
  \node[fd2 lab,left] at (-5.15,0) {$\bar B^0$};
  \node[fd2 lab,right] at (5.15,0) {$D^+$};
  \node[fd2 lab,anchor=south west] at (2.78,4.45) {$\pi^-(p)$};
  \node[fd2 lab] at (-1.9,0.2) {$b$};
  \node[fd2 lab] at (1.9,0.2) {$c$};
  \node[fd2 lab] at (-1.9,-0.72) {$\bar d$};
  \path[mom={$xp$},mom sep=0.25cm,mom length=0.5cm,mom pos=0.8] (W) -- (dq);
  \path[mom'={$(1-x)p$},mom sep=0.25cm,mom length=0.5cm,mom pos=0.84] (W) -- (uq);
  \node[fd2 lab,anchor=east] at (1.28,3.42) {$\phi_\pi(x)$};
  \node[fd2 note,anchor=north west] at (0.14,0.36) {$T(x,\mu)$};
  \node[fd2 note,anchor=south west,align=left] at (3.7,0.98) {软胶子在领头幂次\\退耦（色透明性）};
  % ---- the soft B -> D transition ----
  \draw[fd2 bracket] (-4.02,-1.05) -- (-4.02,-1.15) -- (4.02,-1.15) -- (4.02,-1.05);
  \node[fd2 note,below=1pt] at (0,-1.15) {软的 $B\to D$ 跃迁（含旁观者）：形状因子 $F^{B\to D}(0)$};
  \node[fd2 lab] at (0,-2.2) {$\langle D^+\pi^-|O_{\mathbf 0}|\bar B^0\rangle=N\,F^{B\to D}(0)\displaystyle\int_0^1\!\dd x\,T(x,\mu)\,\phi_\pi(x,\mu)\,\big[1+\order{\Lambda/m_b}\big]$};
\end{tikzpicture}
\caption{$\bar B^0\to D^+\pi^-$ 的因子化（\cref{thm:heavy-bdpi}，全阶论证 [A]）。黑方块为弱作用四夸克算符 $O_{\mathbf 0}=[\bar c\gamma^\mu P_Lb][\bar d\gamma_\mu P_Lu]$ 的顶点，粗线为重夸克 $b$、$c$，双线为介子；细线上的箭头为费米子数流向（反夸克 $\bar d$、$\bar u$ 的箭头逆着动量）。
旁观者 $\bar d$ 自 $\bar B^0$ 进入 $D^+$，这一软的跃迁由形状因子 $F^{B\to D}(0)$ 描述；发射的 $\pi^-$ 由弱顶点直接产生的共线 $d\bar u$ 对构成（蓝色团，分布振幅 $\phi_\pi(x)$，$x$ 为 $d$ 夸克的动量分数），它是快速运动的色单态小偶极子：
软胶子与 $d$、$\bar u$ 的程函耦合在领头幂次相消（打叉的绿色胶子，色透明性），硬胶子交换归入短程系数 $T(x,\mu)$。
$N=\frac{\ii}{2}m_BE_\pi f_\pi$；修正被 $\Lambda/m_b$、$\Lambda/m_c$、$\Lambda/E_\pi$ 压低（VF-BQ-6）；$B^-\to D^0\pi^-$ 同理。证明还用到分布振幅端点压低的假设（见正文）。}
\label{fig:heavy-bdpi}
\end{figure}
%%% END fig:heavy-bdpi
\begin{example}[朴素因子化与归一化 $N$]\label{ex:heavy-naive}
$\mathcal H_W=\frac{4G_F}{\sqrt2}V_{cb}V_{ud}^*[C_{\mathbf 0}O_{\mathbf 0}+C_{\mathbf 8}O_{\mathbf 8}]$，其中\cite[式~(2)]{BauerPirjolStewart2001}
\begin{equation}\label{eq:heavy-O08}
O_{\mathbf 0}=\big[\bar c\gamma^\mu P_Lb\big]\big[\bar d\gamma_\mu P_Lu\big],\qquad
O_{\mathbf 8}=\big[\bar c\gamma^\mu P_Lt^ab\big]\big[\bar d\gamma_\mu P_Lt^au\big],\qquad P_L=\tfrac12(1-\gamma_5).
\end{equation}
% VF-BQ-6
在 $\bar B$ 静止系中取无质量的 $\pi^-$ 沿 $+z$，$p_\pi=\sqrt2E_\pi n_A$，$E_\pi=\frac{m_B^2-m_D^2}{2m_B}=2.309$\,GeV（$m_B=5.280$\,GeV，$m_D=1.870$\,GeV\cite{ParticleDataGroup:2024cfk}）。真空插入把矩阵元写成两个流的矩阵元之积。由\cref{def:excl-da}的局域极限 $\langle0|\bar d\gamma^\mu\gamma_5u|\pi^+(p)\rangle=\ii f_\pi p^\mu$、厄米共轭与同位旋，$\langle\pi^-|\bar d\gamma_\mu\gamma_5u|0\rangle=-\ii f_\pi p_{\pi\mu}$，矢量流部分由宇称为零，故 $\langle\pi^-|\bar d\gamma_\mu P_Lu|0\rangle=\frac{\ii}{2}f_\pi p_{\pi\mu}$。$\bar B\to D$ 只有矢量流贡献：按 $\langle D(p')|\bar c\gamma^\mu b|\bar B(p)\rangle=F_+[(p+p')^\mu-\frac{m_B^2-m_D^2}{q^2}q^\mu]+F_0\frac{m_B^2-m_D^2}{q^2}q^\mu$ 分解，与 $q=p_\pi$ 收缩时 $F_+$ 项消去，剩 $\frac12(m_B^2-m_D^2)F_0(0)$，$F_0(0)=F_+(0)\equiv F^{B\to D}(0)$。于是
\begin{equation}\label{eq:heavy-naive}
\langle D^+\pi^-|O_{\mathbf 0}|\bar B^0\rangle\big|_{\rm naive}=\frac{\ii}{2}f_\pi\cdot\frac12(m_B^2-m_D^2)F^{B\to D}(0)=\frac{\ii}{2}m_BE_\pi f_\pi\,F^{B\to D}(0)\equiv N\,F^{B\to D}(0),
\end{equation}
与\cite[式~(8)]{BauerPirjolStewart2001}一致；用 $(V-A)\otimes(V-A)$ 算符时 $N$ 大 $4$ 倍，$O_{\mathbf 8}$ 的真空插入 $\propto\Tr t^a=0$。重夸克极限下 $F^{B\to D}(0)=\frac{m_B+m_D}{2\sqrt{m_Bm_D}}\xi(w)$（$\xi$ 为 Isgur--Wise 函数，$w=\frac{m_B^2+m_D^2}{2m_Bm_D}=1.589$，前置因子 $1.138$；脚本 \texttt{s1\_\allowbreak bdpi.py}）。
\uses{def:excl-da}
\end{example}

\paragraph{色透明性。}
软胶子（$k\sim\Lqcd$）与出射夸克 $p_1=xp_\pi+l$、反夸克 $p_2=\bar xp_\pi-l$ 耦合（$\bar x=1-x$，$|l|\sim\Lqcd$，$p_\pi\cdot l=\order{l^2}$）。色单态的 $q\bar q$ 对的两条线给出同一色矩阵而符号相反，展开到 $l$ 的一阶，
\begin{equation}\label{eq:heavy-dipole}
J^\mu\propto t^a\Big[\frac{p_1^\mu}{p_1\cdot k}-\frac{p_2^\mu}{p_2\cdot k}\Big]
=t^a\Big(\frac1x+\frac1{\bar x}\Big)\frac{1}{p_\pi\cdot k}\Big[l^\mu-p_\pi^\mu\,\frac{l\cdot k}{p_\pi\cdot k}\Big]+\order{l^2}.
\end{equation}
$l=0$ 时 $J^\mu\equiv0$（程函因子只依赖于方向）；余项是守恒的色偶极子（$k_\mu J^\mu=0$），相对单条线为 $\order{\Lqcd/(x\bar xm_b)}$——\emph{端点 $x\bar x\sim\Lqcd/m_b$ 除外}，这就是需要端点压低假设的原因（脚本 \texttt{s1\_\allowbreak bdpi.py}）。

\begin{theorem}[$B\to D\pi$ 的全阶因子化 \tagA]\label{thm:heavy-bdpi}
设 $Q\in\{m_b,m_c,E_\pi\}$ 同阶且远大于 $\Lqcd$（$m_c/m_b$ 固定），$\lambda=\Lqcd/Q$，并设（端点压低假设）介子波函数的端点组态——$\pi$ 中一个价夸克带走几乎全部能量的组态，以及 $B$、$D$ 中旁观者动量 $\gg\Lqcd$ 的组态——被 $\lambda^{a}$（某个 $a>0$）压低。则对 $\bar B^0\to D^+\pi^-$ 与 $B^-\to D^0\pi^-$，在 $\lambda$ 的领头幂次、对 $\as$ 的所有阶，
\begin{equation}\label{eq:heavy-bdpi}
\langle D\pi|O|B\rangle=N\,F^{B\to D}(0)\int_0^1\dd x\;T_O(x,\mu)\,\phi_\pi(x,\mu)\,\big[1+\order{\lambda}\big],\qquad N=\frac{\ii}{2}\,m_BE_\pi f_\pi ,
\end{equation}
% VF-BQ-6
其中 $O=O_{\mathbf 0}$ 或 $O_{\mathbf 8}$（\cref{eq:heavy-O08}）；$F^{B\to D}(0)$ 是 $q^2=0$ 处的形状因子（取 $F^{B\to D}(m_\pi^2)$ 只差幂次压低的项）；$\phi_\pi$ 是\cref{def:excl-da}的分布振幅，$x$ 为 $\pi^-$ 中 $d$ 夸克的动量分数；$T_O$ 是可微扰计算的短程系数（依赖于 $\mu$、$m_c/m_b$ 与匹配标度），$T_{O_{\mathbf 0}}=1+\order{\as}$，$T_{O_{\mathbf 8}}=\order{\as}$。
非微扰地位：\cref{asm:intro-bridge}（MA1、MA2）。
% errtype: E-S
\uses{def:excl-da,thm:scet-decoupling,asm:intro-bridge}
\end{theorem}

\begin{proofsketch}
按\cite{BauerPirjolStewart2001}，在 SCET 与 HQET 中进行。
(1)~动量 $\gtrsim Q$ 的自由度积出后进入 Wilson 系数；剩下 $B$、$D$ 中的 HQET 重夸克与软夸克、软胶子（$\sim Q\lambda$），$\pi$ 的共线夸克 $\xi$（$\sla n_A\xi=0$）与共线胶子（$(p^+,p^-,p_T)\sim Q(1,\lambda^2,\lambda)$），以及超软胶子（$\sim Q\lambda^2$）。它们穷尽领头区域是 SCET 的第一个前提（\cref{rem:scet-assumptions}；此步只给出 E-S 型陈述）。
(2)~共线胶子接到重夸克上使之离壳，积出后求和为沿 $n_B$ 的 Wilson 线 $W$（因 $\sla vh_v=h_v$ 而与速度无关）：$O_{\mathbf 0,\mathbf 8}$ 匹配到 $Q_{\mathbf 0}=[\bar h_{v'}\Gamma_hh_v][\bar\xi\Gamma_\ell\xi]$、$Q_{\mathbf 8}=[\bar h_{v'}\Gamma_h(W^\dagger t^aW)h_v][\bar\xi\Gamma_\ell t^a\xi]$；对色单态 $W^\dagger W=1$，即来自 $b$ 与 $c$ 的共线胶子耦合对所有阶相消。因 $\sla n_A\xi=0$，旁观者不能在 $\order{\lambda^0}$ 变成共线夸克。
(3)~超软胶子与重夸克、软模式的耦合至少被一个 $\lambda$ 压低，与共线场的领头耦合由 BPS 场重定义消去（\cref{thm:scet-decoupling}，[R, $R_{\mathrm{def}}$]）。软胶子在 $\order{\lambda^0}$ 只经 $n_A\cdot A_s$ 与共线粒子耦合，软动量只以 $n_A\cdot k$ 进入共线传播子（$(p_c+k)^2\simeq2p_c^+\,n_A\cdot k$）；由 Ward 恒等式，软胶子可提出成沿 $n_A$ 的软 Wilson 线 $S$。对 $Q_{\mathbf 0}$，$S^\dagger S=1$（\cref{eq:heavy-dipole}的全阶形式）；对 $Q_{\mathbf 8}$，软交换的色结构 $M$ 的伴随指标全部缩并，$\Tr[MW^\dagger t^aW]=M\,\Tr t^a=0$\cite[式~(7)]{BauerPirjolStewart2001}。
(4)~共线规范不变性把硬修正限制为 $[\bar h_{v'}\Gamma_hh_v][\bar\xi W\,C(\bar{\mathcal P}_+)\,\Gamma_\ell W^\dagger\xi]$，$\bar{\mathcal P}_+$ 给出夸克与反夸克大动量之差；$\pi$ 的矩阵元是 $\phi_\pi$ 的 Fourier 变换，故得卷积，$T(x,\mu)=C(\mu,2(2x-1)E_\pi)$，树图时回到\cref{eq:heavy-naive}。
(5)~端点压低假设使弱湮灭、硬旁观者以及旁观者进入 $\pi$ 的拓扑被幂次压低（例如硬旁观者图 $\sim G_F\as m_b\Lqcd^2$，领头振幅 $\sim G_Fm_b^2\Lqcd$\cite[式~(26)、(28)]{BenekeEtAl2000}）。证明不需要 $\pi$ 由 $q\bar q$ Fock 态主导，也不需要令横动量为零\cite{BauerPirjolStewart2001}；$T$ 的两圈结果见\cite{Huber:2016xod}。
\end{proofsketch}

\paragraph{地位与缺口}
全阶论证（[A]）。最弱的一步是 SCET 的前提（\cref{rem:scet-assumptions}，其与本书编号假设的对应见该注）：模式穷尽领头区域——强子层次上即\cref{asm:intro-bridge}(Br1)——与零箱减除的自洽；BPS 的模式清单不含 Glauber 模式，其不出现属于该注的前提~(3)。这些步骤只给出 E-S 型陈述。另一输入是端点压低假设，它是关于强子波函数的非微扰陈述，BPS 明确列为假设，没有证明。BBNS 在两圈显式证明了红外奇异性的相消，并写明“the following does not accomplish an all-order ‘proof’”（以下论证并不构成全阶“证明”）\cite[第~5~节]{BenekeEtAl2000}。没有已知的领头幂次反例。升级途径：对本过程以积分形式证明模式完备，并控制端点组态。

% ---------------------------------------------------------------------
\section{重到轻形状因子与无粲 $B$ 衰变的 QCD 因子化}\label{sec:heavy-qcdf}

两个末态介子都很轻时（$\bar B\to\pi\pi$、$\pi K$ 等），接收旁观反夸克的介子 $M_1$ 也是高能的：旁观者要么以趋于零的动量分数进入 $M_1$（软重叠），要么经虚度 $\sim m_b\Lqcd$ 的硬共线胶子交换获得大动量（硬旁观者散射，\cref{fig:heavy-qcdf}(b)），于是出现中间标度 $\sqrt{m_b\Lqcd}$。

%%% BEGIN fig:heavy-qcdf
\begin{figure}[htbp]
\centering
\begin{tikzpicture}[fd,
  fd2 heavy/.style={draw,solid,line width=1.3pt,fd arrow=#1},
  fd2 heavy/.default=0.5,
  fd2 meson/.style={bound state,line cap=butt},
  fd2 weak/.style={rectangle,fill=black,draw=none,inner sep=0pt,minimum size=5.6pt},
  fd2 halo/.style={preaction={draw=white,line width=3.4pt}},
  fd2 endpoint/.style={draw=red!80!black,line width=3.2pt,opacity=0.45,line cap=butt},
  fd2 bracket/.style={draw=gray!55!black,line width=0.5pt},
  fd2 lab/.style={font=\small},
  fd2 note/.style={font=\footnotesize}]
  % ------------- common skeleton: B (left), M1 (right), M2 (emitted, top) -------------
  \foreach \panel/\xs in {a/0, b/8.0} {
  \begin{scope}[xshift=\xs cm]
    \coordinate (W) at (-0.5,0.42);
    \coordinate (q2) at (0.352,2.160);   % quark of M2 (inside the blob)
    \coordinate (a2) at (0.698,1.960);   % antiquark of M2 (inside the blob)
    \draw[fd2 meson] (-3.3,0) -- (-2.75,0);
    \draw[fd2 meson] (2.75,0) -- (3.3,0);
    \draw[fd2 meson] (1.08,3.03) -- (1.43,3.64);
    \if a\panel
      % hard gluon of a vertex correction: M2 antiquark -- light quark q
      \coordinate (g1) at (1.2,0.42);  \coordinate (g2) at ($(W)!0.58!(a2)$);
      \draw[gluon] (g1) -- (g2);
      \draw[fd2 heavy=0.5] (-2.3,0.42) -- (W);
      \draw[quark=0.78] (W) -- (2.3,0.42);
      \draw[quark] (2.3,-0.42) -- (-2.3,-0.42);
      \node[vertex] at (g1) {};  \node[vertex] at (g2) {};
    \else
      % hard-collinear gluon from the spectator to the emitted pair
      \coordinate (g1) at (0.85,-0.42);  \coordinate (g2) at ($(W)!0.52!(q2)$);
      \draw[fd2 endpoint] (0.9,-0.42) -- (2.2,-0.42);
      \draw[gluon] (g1) -- (g2);
      \draw[fd2 heavy=0.5] (-2.3,0.42) -- (W);
      \draw[quark=0.8,fd2 halo] (W) -- (2.3,0.42);
      \draw[quark=0.55] (2.3,-0.42) -- (g1);
      \draw[quark=0.45] (g1) -- (-2.3,-0.42);
      \node[vertex] at (g1) {};  \node[vertex] at (g2) {};
    \fi
    \draw[quark=0.33] (W) -- (q2);
    \draw[quark=0.6] (a2) -- (W);
    \node[hadron blob,ellipse,minimum width=0.66cm,minimum height=1.56cm,inner sep=0pt] at (-2.45,0) {};
    \node[jetB blob,ellipse,minimum width=0.66cm,minimum height=1.56cm,inner sep=0pt] at (2.45,0) {};
    \node[jet blob,ellipse,minimum width=1.35cm,minimum height=0.62cm,inner sep=0pt,rotate=60] at (0.75,2.45) {};
    \node[fd2 weak] at (W) {};
    \node[fd2 lab,left] at (-3.3,0) {$\bar B$};
    \node[fd2 lab,right] at (3.3,0) {$M_1$};
    \node[fd2 lab,anchor=south west] at (1.36,3.55) {$M_2$};
    \node[fd2 lab] at (-1.45,0.66) {$b$};
    \node[fd2 lab] at (1.5,0.66) {$q$};
    \node[fd2 lab,anchor=east] at (0.3,2.72) {$\phi_{M_2}(x)$};
  \end{scope}}
  % ------------- panel (a): labels -------------
  \node[fd2 lab] at (-1.45,-0.66) {$\bar q_s$};
  \node[fd2 note,anchor=west] at (1.05,1.25) {$T^{\rm I}$：硬胶子};
  \draw[fd2 bracket] (-2.82,-1.0) -- (-2.82,-1.1) -- (2.82,-1.1) -- (2.82,-1.0);
  \node[fd2 note,below=1pt] at (0,-1.1) {软的 $B\to M_1$ 跃迁：$F^{B\to M_1}(0)$};
  \node[fd2 lab,align=center] at (0,-2.3) {(a)\ 顶点项\\[2pt]$F^{B\to M_1}(0)\displaystyle\int\!\dd x\,T^{\rm I}(x)\,\phi_{M_2}(x)$};
  % ------------- panel (b): labels -------------
  \begin{scope}[xshift=8.0cm]
    \path[mom={$\xi p_B$},mom sep=0.2cm,mom length=0.5cm,mom pos=0.4] (-2.2,-0.42) -- (0.85,-0.42);
    \path[mom={$\bar y\,p_1$},mom sep=0.2cm,mom length=0.45cm,mom pos=0.52] (0.85,-0.42) -- (2.2,-0.42);
    \node[fd2 note,anchor=south east] at (-0.12,1.06) {$T^{\rm II}$};
    \node[fd2 note,anchor=west] at (0.56,1.02) {$k^2\sim m_b\Lambda$};
    \node[fd2 note] at (-2.45,-1.1) {$\phi_B(\xi)$};
    \node[fd2 note] at (2.6,-1.12) {$\phi_{M_1}(y)$};
    \node[fd2 note,text=red!70!black] at (1.05,-0.74) {端点区 $\bar y\to0$};
    \node[fd2 lab,align=center] at (0,-2.3) {(b)\ 硬旁观者散射\\[2pt]$\displaystyle\int\!\dd\xi\,\dd x\,\dd y\,T^{\rm II}(\xi,x,y)\,\phi_B(\xi)\,\phi_{M_1}(y)\,\phi_{M_2}(x)$};
  \end{scope}
\end{tikzpicture}
\caption{无粲两体衰变 $\bar B\to M_1M_2$ 的 QCD 因子化（BBNS；\cref{prop:heavy-qcdf-lp}，部分 [P]），各只画一个代表图。黑方块为弱作用四夸克算符的顶点，粗线为 $b$ 夸克，箭头为费米子数流向；
$M_2$（蓝）为由弱顶点发射的介子，$M_1$（橙）为接收旁观反夸克 $\bar q_s$ 的介子。(a)~形状因子项：硬胶子连接发射对与 $b\to q$ 线（顶点修正），软的跃迁由 $F^{B\to M_1}(0)$ 描述。
(b)~硬旁观者散射：旁观者（动量 $\xi p_B$，$\xi\sim\Lambda/m_b$）放出一个硬共线胶子（虚度 $\sim m_b\Lambda$），以动量分数 $\bar y=1-y$ 进入 $M_1$。
红色标出端点区 $\bar y\to0$（旁观反夸克以趋于零的动量分数进入 $M_1$）：扭度 3 的硬旁观者项在此对数发散，$X_H=\int_0^1\frac{\dd y}{1-y}\Phi_p(y)$；弱湮灭项（未画出）有同类端点发散 $X_A=\int_0^1\frac{\dd y}{y}$（VF-BQ-3）。
这些幂次修正的标准共线因子化不成立（\cref{prop:heavy-qcdf-endpoint}，[V]），$X_{A,H}$ 的参数化只是模型。}
\label{fig:heavy-qcdf}
\end{figure}
%%% END fig:heavy-qcdf
\begin{proposition}[重到轻形状因子：软重叠与硬旁观者 \tagA]\label{prop:heavy-btopi-ff}
设 $\bar B\to\pi$ 跃迁中 $\pi$ 在 $\bar B$ 静止系中的能量 $E=\frac{m_B^2-q^2}{2m_B}\sim m_b/2\gg\Lqcd$（大反冲，略去 $m_\pi$）。则在 $\Lqcd/m_b$ 的领头幂次、对 $\as$ 的所有阶，三个形状因子 $F_i$（$i=+,0,T$，定义同\cite{Beneke:2003pa}）可写成
\begin{equation}\label{eq:heavy-btopi-ff}
F_i^{B\to\pi}(E)=C_i(E,\mu)\,\xi_\pi(E,\mu)+\int_0^\infty\!\dd\omega\int_0^1\!\dd u\;T_i(E;\omega,u,\mu)\,\phi_B^+(\omega,\mu)\,\phi_\pi(u,\mu),
\end{equation}
其中 (a)~$C_i$、$T_i$ 由标度 $m_b$ 与 $\sqrt{m_b\Lqcd}$ 上的短程物理决定，可按 $\as$ 展开，$C_i$ 的树图值为运动学因子，$T_i=\order{\as}$ 是硬系数与硬共线喷注函数的卷积；(b)~“软重叠”（soft overlap）形状因子 $\xi_\pi$ 与 $i$ 无关，在 SCET 中定义为同时包含硬共线、共线与软效应的矩阵元；(c)~对 $\omega$、$u$ 的积分收敛。$\phi_B^+$ 是 $B$ 介子的领头光锥分布振幅（$\omega$ 为旁观者光锥动量，$\int_0^\infty\frac{\dd\omega}{\omega}\phi_B^+=\lambda_B^{-1}$\cite{Grozin:1996pq}，演化见\cite{Lange:2003ff}），$\phi_\pi$ 见\cref{def:excl-da}。两项都 $\sim(\Lqcd/m_b)^{3/2}$。
非微扰地位：\cref{asm:intro-bridge}（MA1、MA2）。
% errtype: E-S
\uses{def:excl-da,thm:scet-decoupling,asm:intro-bridge}
\end{proposition}

\begin{proofsketch}
(1)~两步匹配 QCD $\to$ SCET$_{\rm I}$ $\to$ SCET$_{\rm II}$，依次积出硬模式与硬共线模式；在 SCET$_{\rm I}$ 中 $B$ 的软组分相对于硬共线模式是超软的，经 BPS 场重定义退耦（\cref{thm:scet-decoupling}）。(2)~只含一个硬共线夸克场 $\bar\xi W$ 的重到轻流，因 $\sla vh_v=h_v$、$\sla n_A\xi=0$，对不同 Dirac 结构只差运动学因子，给出 $\xi_\pi$ 的普适性；多含一个硬共线胶子场的流在积出硬共线模式后化为喷注函数与 $\phi_B^+$、$\phi_\pi$ 的卷积（E-S 型陈述）。(3)~前一类矩阵元写成卷积会在端点发散；Beneke 与 Feldmann 把这些效应全部吸收进 $\xi_\pi$ 的定义，证明了\cref{eq:heavy-btopi-ff}对微扰论所有阶成立且卷积收敛\cite{Beneke:2003pa}；BPS 用另一种标度分离得到无端点奇异性的硬核，并表明通常归于波函数“尾部”的软效应与硬旁观者项同阶\cite{Bauer:2002aj}（另见\cite{Hill:2002vw,Beneke:2005gs}）。
\end{proofsketch}

\paragraph{地位与缺口}
全阶论证（[A]），前提是\cref{rem:scet-assumptions}的 SCET 前提（模式完备与零箱减除，E-S 型陈述）。被因子化的只是“软重叠 + 硬旁观者”的分解：$\xi_\pi$ 不能再写成分布振幅的卷积，否则端点发散（同类现象见\cref{prop:heavy-qcdf-endpoint}与\cref{prop:excl-endpoint}，[P]）；对 $b$、$c$ 都按非相对论处理的 $B_c\to\eta_c$，软重叠中双对数的全阶结构已经分析，由一组尚无闭式解的耦合积分方程描述\cite{Bell:2024bxg}。升级途径：对 $B\to\pi$ 以积分形式证明模式完备。

\paragraph{记号与幂次计数。}
以 $b\to s$ 为例，$\mathcal H_{\rm eff}=\frac{G_F}{\sqrt2}\sum_{p=u,c}\lambda_p(C_1Q_1^p+C_2Q_2^p+\cdots)$，$\lambda_p=V_{pb}V_{ps}^*$，$Q_1^p=(\bar pb)_{V-A}(\bar sp)_{V-A}$，$Q_2^p=(\bar p_ib_j)_{V-A}(\bar s_jp_i)_{V-A}$，省略号为企鹅算符；NDR 方案中 $\mu=m_b=4.2$\,GeV 处 $C_1=1.081$、$C_2=-0.190$\cite[表~1]{BenekeEtAl2001}。领头扭度分布振幅 $\Phi_M(x)=6x\bar x[1+\sum_{n\ge1}\alpha_n^MC_n^{(3/2)}(2x-1)]$（$\alpha_n^M$ 即\cref{def:excl-da}中的 $a_n$），扭度 3 振幅取 $\Phi_p=1$，$r_\chi^M=2\mu_M/m_b$（如 $\mu_\pi=m_\pi^2/(m_u+m_d)$），$\int_0^1\frac{\dd\xi}{\xi}\Phi_B=\frac{m_B}{\lambda_B}$。$M_1$ 接收旁观者，$M_2$ 由弱顶点发射；振幅写成朴素因子化矩阵元 $A_{M_1M_2}=\ii\frac{G_F}{\sqrt2}m_B^2F_0^{B\to M_1}(0)f_{M_2}$\cite[式~(11)]{BenekeNeubert2003}乘以系数 $a_i$。由 $F^{B\to\pi}(0)\sim(\Lqcd/m_b)^{3/2}$、$f_B\sim\Lqcd^{3/2}m_b^{-1/2}$、$f_\pi\sim\lambda_B\sim\Lqcd$，形状因子项 $\sim m_b^2F^{B\to\pi}(0)f_\pi$ 与硬旁观者项 $\sim f_Bf_{M_1}f_{M_2}m_b/\lambda_B$ 同为 $\Lqcd^{5/2}m_b^{1/2}$。旁观者（动量 $\xi p_B$，$\xi\sim\Lqcd/m_b$）以分数 $\bar y=1-y$ 进入 $M_1$ 时，交换胶子的动量 $k=\xi p_B-\bar yp_1$ 满足（用 $2p_B\cdot p_1=m_B^2$）
\begin{equation}\label{eq:heavy-hc-gluon}
k^2=\xi^2m_B^2-\bar y\,\xi\,m_B^2\sim-m_b\Lqcd\qquad(\bar y=\order{1}),
\end{equation}
是硬共线的；BBNS 因此在 $\mu_h=\sqrt{\Lambda_h\mu}$（$\Lambda_h=0.5$\,GeV）处取硬旁观者项的 $\as$\cite{BenekeEtAl2001}。

\begin{proposition}[无粲两体衰变的领头幂次 QCD 因子化 \tagP]\label{prop:heavy-qcdf-lp}
设 $M_1$、$M_2$ 为味非单态的轻赝标介子。BBNS 因子化公式
\begin{equation}\label{eq:heavy-qcdf}
\begin{split}
\langle M_1M_2|Q_i|\bar B\rangle={}&F^{B\to M_1}(0)\int_0^1\!\dd x\,T^{\rm I}_i(x)\,f_{M_2}\Phi_{M_2}(x)\\
&+\int_0^1\!\dd\xi\,\dd x\,\dd y\;T^{\rm II}_i(\xi,x,y)\,f_B\Phi_B(\xi)\,f_{M_1}\Phi_{M_1}(y)\,f_{M_2}\Phi_{M_2}(x)
\end{split}
\end{equation}
（旁观者可进入任一介子时再加 $M_1\leftrightarrow M_2$ 的项；归一化使树图的 $T^{\rm I}_i$ 恰为朴素因子化，$T^{\rm II}_i=\order{\as}$）作为对 $\as$ 所有阶的领头幂次陈述没有公认的证明；本命题断言它在有限阶已被确立的内容。
(a)~（单圈）在 $\order{\as}$，“非因子化”的顶点修正与旁观者图中的软、共线发散在领头幂次相消，$F^{B\to M_1}$ 按\cref{prop:heavy-btopi-ff}理解；顶点修正给出
\begin{gather}
a_{1,\rm I}=C_1+\frac{C_2}{N_c}\Big[1+\frac{C_F\as}{4\pi}V_{M_2}\Big],\qquad a_{2,\rm I}=C_2+\frac{C_1}{N_c}\Big[1+\frac{C_F\as}{4\pi}V_{M_2}\Big],\label{eq:heavy-ai}\\
V_M=12\ln\frac{m_b}{\mu}-18+\int_0^1\!\dd x\,g(x)\,\Phi_M(x)\quad(\text{NDR}),\label{eq:heavy-VM}\\
\int_0^1\!\dd x\,g(x)\Phi_M(x)=-\frac12-3\ii\pi+\Big(\frac{11}{2}-3\ii\pi\Big)\alpha_1^M-\frac{21}{20}\alpha_2^M+\cdots,\label{eq:heavy-gegen}
\end{gather}
% VF-BQ-1, VF-BQ-2
其中 $g(x)=3\big(\frac{1-2x}{1-x}\ln x-\ii\pi\big)+\big[2\,\mathrm{Li}_2(x)-\ln^2x+\frac{2\ln x}{1-x}-(3+2\ii\pi)\ln x-(x\leftrightarrow1-x)\big]$（$x$ 的约定同\cite{BenekeEtAl2001}），系数 $12$ 由 $a_{1,2}$ 与 $\mu$ 无关唯一确定（对任意 $N_c$）。
(b)~（高阶）$T^{\rm I}$ 的红外有限性已显式验证到两圈（NNLO），$T^{\rm II}$ 的到其单圈修正（$\order{\as^2}$）。
非微扰地位：\cref{asm:intro-bridge}（MA1、MA2）。
% errtype: E-S
\uses{prop:heavy-btopi-ff,def:excl-da,asm:intro-bridge}
\end{proposition}

\begin{proofsketch}
(a)~是 BBNS 的单圈计算\cite[式~(46)--(48)]{BenekeEtAl2001}。系数 $12$ 可独立导出：设 $\frac{\dd\bm C}{\dd\ln\mu}=\frac{\as}{4\pi}\gamma^{\tp}\bm C$，$\gamma=\begin{psmallmatrix}-6/N_c&6\\6&-6/N_c\end{psmallmatrix}$（$C_\pm=C_1\pm C_2$ 的本征值 $\pm6(N_c\mp1)/N_c$；$\mu=m_b$ 处 $C_+=0.891<1<C_-=1.271$，与其符号相符），则
\begin{equation}\label{eq:heavy-a1-rg}
\frac{\dd}{\dd\ln\mu}\Big(C_1+\frac{C_2}{N_c}\Big)=\frac{\as}{4\pi}\Big[\Big(-\frac{6}{N_c}C_1+6C_2\Big)+\frac1{N_c}\Big(6C_1-\frac6{N_c}C_2\Big)\Big]=\frac{\as}{4\pi}\,\frac{12C_F}{N_c}\,C_2 ,
\end{equation}
恰被 $\frac{C_2}{N_c}\frac{C_F\as}{4\pi}12\ln\frac{m_b}{\mu}$ 的导数抵消（准确到 $\order{\as}$）；$a_{2,\rm I}$ 同理，常数 $-18$ 依赖于方案。\cref{eq:heavy-gegen}的系数由数值积分重现到 $10^{-20}$（脚本 \texttt{s2\_\allowbreak qcdf.py}）。(b)~两圈顶点\cite{Bell:2007tv,Bell:2009nk,Beneke:2009ek}与企鹅振幅\cite{Bell:2015koa,Bell:2020qus}，单圈旁观者散射\cite{Beneke:2005vv,Beneke:2006mk,Kivel:2006xc,Pilipp:2007mg,Beneke:2005gs}：这些硬散射核都没有残留的软或共线发散。
\end{proofsketch}

\begin{example}[顶点修正与微扰强相位]\label{ex:heavy-vertex}
取渐近分布振幅与 $\mu=m_b$，$V_\pi=-18.5-3\ii\pi$；以 $\as(m_b)=0.223$（由 $\as(M_Z)=0.1180$ 两圈跑动）得 $\frac{C_F\as}{4\pi}=0.0237$，$a_{1,\rm I}\approx1.045+0.014\,\ii$，$a_{2,\rm I}\approx0.012-0.080\,\ii$。领头幂次的强相位是 $\order{\as}$ 的微扰量，来自 $g(x)$ 中的 $\ii\pi$，即顶点图中轻夸克中间态的在壳部分；色压低的 $a_2$ 中修正与树图值 $C_2+C_1/N_c=0.170$ 同量级，BBNS 指出那里硬旁观者贡献是主要效应，预言相应地不确定\cite[第~4.2~节]{BenekeEtAl2001}（脚本 \texttt{s2\_\allowbreak qcdf.py}）。
\uses{prop:heavy-qcdf-lp}
\end{example}

\paragraph{地位与缺口}
部分（[P]）。一般断言（BBNS 公式在领头幂次对 $\as$ 所有阶成立）没有公认的证明；BBNS 在给出 $B\to D\pi$ 的两圈证明时指出须把它推广到两个轻介子的末态\cite[第~8~节]{BenekeEtAl2000}。SCET 中有领头幂次的全阶推导：Chay 与 Kim 论证振幅对 $\as$ 所有阶因子化\cite{Chay:2003ju}；Bauer、Pirjol、Rothstein、Stewart（BPRS）认为其定理在 $\as(E_M\Lqcd)$ 不可微扰时也成立，并认为粲夸克企鹅（charming penguin）可能在领头阶给出长程效应\cite{Bauer:2004tj}；BBNS 认为二者的公式等价\cite{Beneke:2004bn}，BPRS 不同意其批评\cite{Bauer:2005wb}。本书不作裁断；分歧集中在硬旁观者项能否按 $\as(\sqrt{m_b\Lqcd})$ 展开与粲企鹅的幂次计数（E-S 型陈述）。正因没有共识与公认证明，标签取 [P] 而非 [A]；但说“BBNS 只在 $\order{\as}$ 成立”则低估了现状。升级途径：包含上述两点、被普遍接受的全阶证明。

\begin{proposition}[幂次修正中的端点发散 \tagV]\label{prop:heavy-qcdf-endpoint}
在 $\order{\as}$ 按标准共线因子化计算 $\bar B\to M_1M_2$ 的扭度 3 手征增强硬旁观者项与弱湮灭项（$\Lqcd/m_b$ 幂次修正；$r_\chi^M$ 形式上 $\sim\Lqcd/m_b$，数值上 $\order{1}$）。对端点行为一般的分布振幅（扭度 2 振幅在端点线性趋于零，扭度 3 振幅在端点不为零），尤其对渐近振幅，卷积对数发散，分别含
\begin{equation}\label{eq:heavy-XAXH}
X_H=\int_0^1\frac{\dd y}{1-y}\,\Phi_p(y),\qquad X_A=\int_0^1\frac{\dd y}{y},
\end{equation}
% VF-BQ-3
故这些修正的标准共线因子化不成立。BBNS 的参数化 $X_{A,H}=(1+\varrho_{A,H}\ee^{\ii\varphi_{A,H}})\ln\frac{m_B}{\Lambda_h}$（$\varrho_{A,H}\le1$，相位任意，$\Lambda_h=0.5$\,GeV）是模型，不是因子化定理；例如对两个赝标介子、按 Beneke--Neubert 的规则 $\int_0^1\frac{\dd y}{y}\to X_A$\cite[式~(56)]{BenekeNeubert2003}，
\begin{equation}\label{eq:heavy-annih}
A_1^i\approx\pi\as\Big[18\Big(X_A-4+\frac{\pi^2}{3}\Big)+2r_\chi^2X_A^2\Big],\qquad A_3^f\approx12\pi\as\,r_\chi\big(2X_A^2-X_A\big).
\end{equation}
% VF-BQ-4, VF-BQ-5
（本命题是关于硬散射核与分布振幅之卷积的陈述。）
\uses{prop:heavy-qcdf-lp}
\end{proposition}

\begin{proof}[发散性的显式计算]
\emph{硬旁观者项。}$T^{\rm II}$ 给出（以 $\Phi_{M_2}(x)/\bar x$ 加权并对 $x$ 积分的）$\int_0^1\frac{\dd y}{\bar y}[\Phi_{M_1}(y)+r_\chi^{M_1}\frac{\bar x}{x}\Phi_p(y)]$\cite[式~(57)]{BenekeEtAl2001}。扭度 2 部分有限（渐近振幅给出 $\int_0^16y\,\dd y=3$）；扭度 3 部分为 $r_\chi^{M_1}\langle x^{-1}\rangle_{M_2}X_H$，而 $\int_0^{1-\delta}\frac{\dd y}{1-y}\Phi_p(y)=\Phi_p(1)\ln\frac1\delta+\order{1}$。发散来自旁观者作为软夸克进入 $M_1$ 的区域 $\bar y\to0$，此时由\cref{eq:heavy-hc-gluon} $k^2\to\xi^2m_B^2\sim\Lqcd^2$，交换胶子不再是硬共线的；物理上它在 $\bar y\sim\Lqcd/m_b$ 处被截断，$X_H\sim\ln(m_b/\Lqcd)$，系数可以是复数\cite{BenekeEtAl2001}。
\emph{弱湮灭项。}$A_1^i=\pi\as\int\!\!\int\dd x\,\dd y\,\{\Phi_{M_2}(x)\Phi_{M_1}(y)[\frac{1}{y(1-x\bar y)}+\frac{1}{\bar x^2y}]+r_\chi^{M_1}r_\chi^{M_2}\frac{2}{\bar xy}\}$\cite[式~(62)]{BenekeEtAl2001}。取渐近振幅，第一项有限，$36\int\!\!\int\frac{x\bar x\bar y}{1-x\bar y}=18(\frac{\pi^2}{3}-3)$；第二项 $=18\int_0^1\frac{x\,\dd x}{1-x}$ 与第三项 $=2r_\chi^2\int_0^1\frac{\dd x}{1-x}\int_0^1\frac{\dd y}{y}$ 发散（一般的扭度 2 振幅在 $x\to1$ 处同样给出 $\Phi_{M_2}/\bar x^2\propto1/\bar x$）。把 $\int_0^1\frac{\dd y}{y}$、$\int_0^1\frac{\dd x}{1-x}$ 都换成 $X_A$，用 $\frac{x}{1-x}=\frac1{1-x}-1$，三项为 $18(\frac{\pi^2}{3}-3)$、$18(X_A-1)$、$2r_\chi^2X_A^2$，相加即\cref{eq:heavy-annih}。$A_3^f$ 的两项\cite[式~(62)]{BenekeEtAl2001}在用 $\frac{x(1+\bar x)}{\bar x}=\frac{1}{\bar x}-\bar x$、$\frac{\bar y(1+y)}{y}=\frac1y-y$ 后各给出 $6\pi\as r_\chi(2X_A^2-X_A)$。可见“常数项”只在为发散积分规定数值之后才有意义（脚本 \texttt{s2\_\allowbreak qcdf.py}、\texttt{s8\_\allowbreak review.py}）。
\end{proof}

\paragraph{地位与缺口}
破坏（[V]），证据为显式计算：在标准共线因子化中这些幂次修正的卷积发散。$X_{A,H}$ 只是对未知非微扰贡献的参数化：BBNS 称对所有湮灭项取同一个 $X_A$“is a crude model”（是一个粗糙的模型）\cite[第~3.5~节]{BenekeEtAl2001}，Beneke 与 Neubert 称弱湮灭的处理“is model-dependent”（依赖于模型）\cite[第~2.4~节]{BenekeNeubert2003}。两文把“默认值”印作 $\ln(m_M/\Lambda_h)\approx2.4$，与所用参数化不符，应为 $\ln(m_B/\Lambda_h)=2.36$（附录~Q 第~7~条）。把这些端点发散再因子化成有限的因子化定理，对本过程是猜想性质的（[C]）。其他过程中已有逐例的成功：$h\to\gamma\gamma$（经 $b$ 夸克圈）\cite{Liu:2019oav,Liu:2020wbn}、$gg\to h$（经轻夸克圈）\cite{Liu:2022ajh}、$\bar B\to X_s\gamma$ 的次领头 $O_8$--$O_8$ 贡献\cite{Hurth:2023paz}。对纯弱湮灭衰变，据该文摘要，Lü 等人指出一个此前未计入的领头幂次硬共线贡献，认为它有助于理解卷积中的快度发散\cite{Lu:2022kos}；对纯湮灭道 $\bar B^0\to K^+K^-$，一篇 2024 年的博士论文（未在期刊发表）在树图匹配层次上论证了端点发散在标准贡献与新出现的贡献（含 $B$ 介子尚未研究的四、五粒子软函数）之间相消\cite{Stillger:2024sgy}。色压低的 $\bar B^0\to D^0M^0$ 中，弱湮灭与色压低振幅同阶，没有幂次压低\cite{Mantry:2003uz}。据我们截至 2026 年 10 月的 INSPIRE 检索，对 $\bar B\to M_1M_2$ 的一般弱湮灭项与扭度 3 硬旁观者项尚无经同行评议的全阶再因子化；一般情形见第~21~章的\cref{prop:power-nlp-endpoint}（[P]）与\cref{sec:front-nlp}。解决途径：为本过程建立再因子化——那将是一个形式不同的新定理，不改变标准形式失效这一结论。

% ---------------------------------------------------------------------
\section{单举 $B$ 衰变}\label{sec:heavy-inclusive}

单举衰变率与 $R_{e^+e^-}$、DIS 一样是编时乘积前向矩阵元的吸收部分（如 $\Gamma(\bar B\to X_u\ell\bar\nu)\propto\Imag\langle\bar B|\ii\!\int\!\dd^4x\,T\{J^\dagger(x)J(0)\}|\bar B\rangle$），可按 $1/m_b$ 作算符乘积展开（OPE）；在谱的端点 $p_X^2\sim m_b\Lqcd$，末态是一个喷注，一整串被 $\Lqcd/m_b$ 压低的项与领头项同阶，须求和成形状函数（shape function）。对 $\bar B\to X_s\gamma$，在 $\bar B$ 静止系中写 $p_b=m_bv+k$（剩余动量 $k\sim\Lqcd$），光子 $q=\sqrt2E_\gamma n_B$ 沿 $-z$，$p_X=m_bv+k-q$ 沿 $+z$。由 $p_X^2=m_b^2-2m_bE_\gamma+2m_bk^0-2E_\gamma(k^0+k^3)+k^2$ 与 $2E_\gamma=m_b+\order{\Lqcd}$，
\begin{equation}\label{eq:heavy-pX2}
p_X^2=m_b\big(m_b-2E_\gamma+\omega\big)+\order{\Lqcd^2},\qquad \omega\equiv\frac{n_A\cdot k}{n_A\cdot v}=k^0-k^3 :
\end{equation}
喷注只通过 $\omega$ 一个分量感受 $b$ 夸克的费米运动（脚本 \texttt{s3\_\allowbreak inclusive.py}）。

\begin{theorem}[单举 $B$ 衰变的重夸克展开 \tagA]\label{thm:heavy-inclusive-b}
(a)~（局域 OPE）总半轻与总辐射衰变率——更一般地，对强子末态能量作了足够宽的抹平的单举量，如谱的低阶矩——有按 $\Lqcd/m_b$ 的渐近展开：领头项是自由 $b$ 夸克的衰变率（含其 $\as$ 修正），没有 $\order{\Lqcd/m_b}$ 项，$1/m_b^2$ 项由
\begin{equation}\label{eq:heavy-lambda12}
\lambda_1=\frac{\langle\bar B|\bar h_v(\ii D)^2h_v|\bar B\rangle}{2m_B},\qquad
\lambda_2=\frac{\langle\bar B|\bar h_v\frac g2\sigma_{\mu\nu}G^{\mu\nu}h_v|\bar B\rangle}{6m_B}
\end{equation}
决定（$[\ii D^\mu,\ii D^\nu]=\ii gG^{\mu\nu}$；$m_{B^*}^2-m_B^2=4\lambda_2+\order{\Lqcd^3/m_b}$）。例如以极点质量 $m_b$ 表示，
\begin{equation}\label{eq:heavy-Gamma-bu}
\Gamma(\bar B\to X_u\ell\bar\nu)=\frac{G_F^2|V_{ub}|^2m_b^5}{192\pi^3}\Big[1+\order{\as}+\frac{\lambda_1-9\lambda_2}{2m_b^2}+\order{\frac{\as\Lqcd^2}{m_b^2},\frac{\Lqcd^3}{m_b^3}}\Big].
\end{equation}
(b)~（端点区）当 $m_b-2E_\gamma=\order{\Lqcd}$ 时，在 $\Lqcd/m_b$ 的领头幂次、对 $\as$ 的所有阶，
\begin{equation}\label{eq:heavy-HJS}
\frac{\dd\Gamma(\bar B\to X_s\gamma)}{\dd E_\gamma}=2\Gamma_0\,H(m_b,\mu)\int_{2E_\gamma-m_b}^{\bar\Lambda}\dd\omega\;m_b\,J\big(m_b(m_b-2E_\gamma+\omega),\mu\big)\,S(\omega,\mu),
\end{equation}
其中 $\Gamma_0$ 是 $b\to s\gamma$ 的领头阶宽度，硬函数 $H=1+\order{\as}$，$J(p^2,\mu)=\delta(p^2)+\order{\as}$ 是沿 $n_A$ 的夸克喷注函数，$\bar\Lambda=m_B-m_b$，形状函数为
\begin{equation}\label{eq:heavy-shape}
S(\omega)=\frac{1}{2m_B}\Big\langle\bar B(v)\Big|\,\bar h_v\,\delta\Big(\omega-\frac{n_A\cdot\ii D}{n_A\cdot v}\Big)h_v\Big|\bar B(v)\Big\rangle
\end{equation}
（写成双局域形式时两个 $h_v$ 之间是沿 $n_A$ 的类光 Wilson 线）；$\bar B\to X_u\ell\bar\nu$ 的端点区有用同一形状函数写出的类似因子化。
非微扰地位：\cref{asm:intro-wightman}、\cref{asm:intro-composite}、\cref{asm:intro-af-coeff}（MA0、MA4）；端点（形状函数）区的因子化另需\cref{asm:intro-bridge}（MA1、MA2）。
% errtype: E-S
\uses{thm:dis-ope,thm:scet-decoupling,asm:intro-wightman,asm:intro-composite,asm:intro-af-coeff,asm:intro-bridge}
\end{theorem}

\begin{proofsketch}
(1)~对编时乘积在 HQET 算符上作 OPE；它与 DIS 的 OPE 是同一类短程展开，全阶地位相同（\cref{thm:dis-ope}，[A]）。(2)~领头算符 $\bar b\gamma^\mu b$ 的前向矩阵元由 $b$ 数守恒精确归一；量纲 4 的唯一候选 $\bar h_v\,\ii v\cdot D\,h_v$ 由 HQET 运动方程为零，故没有 $1/m_b$ 修正\cite{ChayGeorgiGrinstein1990,BigiUraltsevVainshtein1992}；量纲 5 的动能与色磁算符给出 $\lambda_{1,2}$。Manohar 与 Wise 对 $b\to u$ 求得 $\Gamma/\Gamma_b=1+5E_b-6K_b-2G_b$，$E_b=K_b+G_b$\cite[式~(5.11)]{ManoharWise1994}；在本书约定下 $K_b=-\frac{\lambda_1}{2m_b^2}$（$b$ 夸克的动能，以 $m_b$ 为单位）、$G_b=-\frac{3\lambda_2}{2m_b^2}$（自旋能量），代入得 $1-K_b+3G_b=1+\frac{\lambda_1-9\lambda_2}{2m_b^2}$，即\cref{eq:heavy-Gamma-bu}；其中 $\frac{\lambda_1}{2m_b^2}$ 就是时间膨胀因子 $\langle m_b/E_b\rangle-1$（$\lambda_1=-\langle\bm p^2\rangle$；脚本 \texttt{s3\_\allowbreak inclusive.py}）。(3)~OPE 只在离开物理割线处可靠，围道在小的强子不变质量处靠近割线，故 (a) 只对总宽度与抹平的量成立（与\cref{ch:ee}类似）。(4)~在端点区，$[\Lqcd/(m_b-2E_\gamma)]^n$ 型的项同阶，求和成形状函数\cite{Neubert1994a,Neubert1994b,BigiEtAl1994}。在 SCET 中 $b\to s\gamma$ 流匹配到 SCET 流，$H$ 是匹配系数的模方；$B$ 中动量 $\sim\Lqcd$ 的组分相对于喷注是超软的，经 BPS 场重定义退耦（\cref{thm:scet-decoupling}），留下沿 $n_A$ 的 Wilson 线，即\cref{eq:heavy-shape}；喷注部分给出 $J$。RG 一致性使 Sudakov 对数可以求和\cite{BauerFlemingLuke2000}；图形层次的同一结构见\cite{KorchemskySterman1994}，喷注函数见\cite{Bauer:2003pi,Becher:2006qw}（此步只给出 E-S 型陈述）。
\end{proofsketch}

\paragraph{地位与缺口}
全阶论证（[A]）。(a) 的地位是微扰 OPE 的地位（\cref{thm:dis-ope}，[A]），(b) 的地位是 SCET 因子化的地位（\cref{rem:scet-assumptions}，模式完备，E-S 型陈述）。用于\emph{固定运动学点}上的强子衰变率需要局域对偶（\cref{conj:ee-duality}(b)，[C]）；它不是定理的前提，而是应用所需的额外输入。以极点质量表示时，$m_b^5$ 带有相对量级 $\Lqcd/m_b$ 的红外重整化子歧义\cite{Bigi:1994em,Beneke:1994sw}；OPE 中没有能吸收它的 $1/m_b$ 项，所以它必须与系数级数中相应的发散相消，改用短距离质量后线性的红外敏感性消失，与没有 $1/m_b$ 项相符\cite{Beneke:1994bc}。升级途径：(a) 随微扰 OPE 升级；(b) 需对端点区证明模式完备。

\begin{example}[端点区的树图检验]\label{ex:heavy-photon}
取 $H=1$、$J=\delta(p^2)$：$m_b\,\delta(m_b(m_b-2E_\gamma+\omega))=\delta(\omega-2E_\gamma+m_b)$，\cref{eq:heavy-HJS}化为 $\frac{\dd\Gamma}{\dd E_\gamma}=2\Gamma_0S(2E_\gamma-m_b)$，光子能量 $E_\gamma=\frac12(m_b+\omega)$ 直接映出 $\omega$ 的分布，积分给出 $\Gamma_0\int S=\Gamma_0$。支撑 $\omega\le\bar\Lambda$\cite{Neubert1994b}把端点从 $m_b/2$ 移到 $m_B/2$。树图层次上 $\int S=1$，$\int\omega S=\langle v\cdot k\rangle=0$（运动方程），$\int\omega^2S=\langle(n_A\cdot k)^2\rangle/(n_A\cdot v)^2=-\frac{\lambda_1}{3}$（由 $\langle k^\mu k^\nu\rangle=\frac{\lambda_1}{3}(g^{\mu\nu}-v^\mu v^\nu)$ 与 $n_A^2=0$），故 $\langle E_\gamma\rangle=m_b/2$，方差 $-\frac{\lambda_1}{12}>0$：费米运动把部分子的 $\delta$ 函数谱展宽成宽度 $\sim\Lqcd$ 的峰。数值上 $m_{B^*}=5.325$\,GeV、$m_B=5.280$\,GeV 给出 $\lambda_2\approx0.12$\,GeV$^2$，\cref{eq:heavy-Gamma-bu}中 $-\frac{9\lambda_2}{2m_b^2}\approx-2.3\%$（$m_b=4.8$\,GeV；脚本 \texttt{s7\_\allowbreak extra.py}）。
\uses{thm:heavy-inclusive-b}
\end{example}

% ---------------------------------------------------------------------
\section{重夸克偶素：NRQCD}\label{sec:heavy-nrqcd}

重夸克偶素有三个标度：$m_Q$、相对动量 $m_Qv$ 与结合能 $m_Qv^2$；势模型与格点给出 $v^2\approx0.3$（$c\bar c$）、$0.1$（$b\bar b$）\cite[第~II.A~节]{BodwinBraatenLepage1995}。湮灭发生在距离 $\sim1/m_Q$ 上，束缚态动力学发生在更大的距离上。非相对论 QCD（NRQCD）积出标度 $m_Q$，按 $v$ 与 $\as(m_Q)$ 双重展开\cite{CaswellLepage1986,BodwinBraatenLepage1995}；衰变的因子化有全阶论证，产生的因子化是猜想（\cref{fig:nrqcd}）。

\begin{definition}[NRQCD 及其幂次计数]\label{def:heavy-nrqcd}
NRQCD 的拉氏量为
\begin{equation}\label{eq:heavy-LNRQCD}
\mathcal L_{\rm NRQCD}=\mathcal L_{\rm light}+\psi^\dagger\Big(\ii D_t+\frac{\bm D^2}{2m_Q}\Big)\psi+\chi^\dagger\Big(\ii D_t-\frac{\bm D^2}{2m_Q}\Big)\chi+\delta\mathcal L_{\rm bilinear}+\delta\mathcal L_{\rm 4f},
\end{equation}
其中 $\mathcal L_{\rm light}$ 是轻夸克与胶子的 QCD 拉氏量（\cref{def:prelim-qcd}），$\psi$（$\chi$）是湮灭重夸克（产生重反夸克）的二分量 Pauli 场，$D_t=\partial_t-\ii gA^0$、$\bm D=\bm\nabla+\ii g\bm A$ 是 $D_\mu=\partial_\mu-\ii gA_\mu$ 的分量（$\bm A$ 为 $A^\mu$ 的空间分量）。$\delta\mathcal L_{\rm bilinear}$ 是 $\frac{1}{8m_Q^3}\psi^\dagger\bm D^4\psi$、Darwin 项、自旋—轨道项与 Fermi 项 $\frac{c_F}{2m_Q}\psi^\dagger\bm\sigma\cdot(\ii\bm D\times\bm D)\psi$ 等双线性修正（及 $\chi$ 的对应项），系数 $=1+\order{\as}$；Fermi 项来自 $(\bm\sigma\cdot\bm D)^2=\bm D^2+\bm\sigma\cdot(\ii\bm D\times\bm D)$，与 $g$ 的符号约定无关。$\delta\mathcal L_{\rm 4f}=\sum_n\frac{f_n(\Lambda)}{m_Q^{d_n-4}}\mathcal O_n$ 为四费米子算符项（$\Lambda$ 为截断标度，$d_n$ 为算符量纲），如 $d_n=6$ 的
\begin{equation}\label{eq:heavy-4f}
\begin{aligned}
&\mathcal O_1({}^1S_0)=\psi^\dagger\chi\,\chi^\dagger\psi,&\qquad&\mathcal O_1({}^3S_1)=\psi^\dagger\bm\sigma\chi\cdot\chi^\dagger\bm\sigma\psi,\\
&\mathcal O_8({}^1S_0)=\psi^\dagger t^a\chi\,\chi^\dagger t^a\psi,&\qquad&\mathcal O_8({}^3S_1)=\psi^\dagger\bm\sigma t^a\chi\cdot\chi^\dagger\bm\sigma t^a\psi,
\end{aligned}
\end{equation}
以及含 $\overleftrightarrow{\bm D}$ 的 $P$ 波算符（$d_n=8$）；电磁算符把中间态投影到 QCD 真空，如 $\mathcal O^{\rm EM}_1({}^3S_1)=\psi^\dagger\bm\sigma\chi|0\rangle\cdot\langle0|\chi^\dagger\bm\sigma\psi$。$f_n$ 的虚部描述 $Q\bar Q$ 在距离 $\sim1/m_Q$ 上的湮灭。\emph{速度标度规则}：设 $v$ 为典型相对速度，$\as(m_Qv)\sim v$，则
\begin{equation}\label{eq:heavy-vscaling}
\psi,\chi\sim(m_Qv)^{3/2},\qquad \bm D\sim m_Qv,\qquad D_t\sim m_Qv^2,\qquad g\bm E\sim m_Q^2v^3,\qquad g\bm B\sim m_Q^2v^4
\end{equation}
% VF-NRQ-1
（$g\bm E$、$g\bm B$ 指 $\ii[D_t,\bm D]$ 与 $\ii\bm D\times\bm D$），故 $\delta\mathcal L_{\rm bilinear}$ 相对领头项为 $\order{v^2}$。BBL 用 $D=\partial+\ii gA$，其 $c_{2,3,4}$ 项写成本书约定下的 $\bm E$、$\bm B$ 时须变号。
\uses{def:prelim-qcd}
\end{definition}

\cref{eq:heavy-vscaling}来自量纲分析与运动方程\cite[第~II.C~节]{BodwinBraatenLepage1995}\cite{Lepage:1992tx}：归一化 $\int\dd^3x\,\psi^\dagger\psi\simeq1$（体积 $(m_Qv)^{-3}$）、Schr\"odinger 方程、Gauss 定律 $\bm D\cdot g\bm E\sim g^2\psi^\dagger t^a\psi$ 与 Amp\`ere 定律 $\bm D\times g\bm B\sim D_tg\bm E+g^2\psi^\dagger\bm D\psi/m_Q$。由此，$S$ 波色单态衰变 $\sim\langle\mathcal O_1\rangle/m_Q^2\sim m_Qv^3$；$P$ 波色单态（量纲 8）$\sim(m_Qv)^5/m_Q^4=m_Qv^5$；$\chi_{QJ}$ 中概率 $\sim v^2$ 的色八重态 ${}^3S_1$ 成分 $\sim v^2\cdot m_Qv^3$。后两者同阶，所以 $P$ 波衰变中色八重态不可缺少（脚本 \texttt{s4\_\allowbreak nrqcd\_lo.py}）。

\begin{theorem}[衰变的 NRQCD 因子化 \tagA]\label{thm:heavy-nrqcd-decay}
设 $H$ 为重夸克偶素（态取非相对论归一化 $\langle H(\bm P')|H(\bm P)\rangle=(2\pi)^3\delta^3(\bm P-\bm P')$，在其静止系中取矩阵元），$m_Q$ 为重夸克的极点质量。则对 $\as(m_Q)$ 的所有阶、对 $v$ 逐阶，
\begin{equation}\label{eq:heavy-nrqcd-decay}
\begin{gathered}
\Gamma(H\to\text{轻强子})=\sum_n\frac{2\Imag f_n(\Lambda)}{m_Q^{d_n-4}}\,\langle H|\mathcal O_n(\Lambda)|H\rangle,\\
\Gamma(H\to\text{电磁末态})=\sum_n\frac{2\Imag f_{{\rm EM},n}(\Lambda)}{m_Q^{d_n-4}}\,\langle H|\mathcal O^{\rm EM}_n(\Lambda)|H\rangle,
\end{gathered}
\end{equation}
求和遍历\cref{def:heavy-nrqcd}的四费米子算符；$\Lambda$ 依赖在系数与矩阵元之间相消；$\Imag f_n$ 由在壳 $Q\bar Q$ 振幅的匹配确定，与 $H$ 无关，可按 $\as(m_Q)$ 展开；到 $v$ 的任一给定阶只有有限个算符贡献。
非微扰地位：\cref{asm:intro-bridge}（MA1、MA2）；非微扰地还需局域对偶（\cref{conj:ee-duality}(b)）。
% errtype: E-S
\uses{def:heavy-nrqcd,asm:intro-bridge}
\end{theorem}

\begin{proofsketch}
(1)~衰变率是态的能量虚部的 $-2$ 倍，湮灭对它的贡献来自 $-\delta\mathcal L_{\rm 4f}$ 的期待值，故 $\Gamma=2\Imag\langle H|\delta\mathcal L_{\rm 4f}|H\rangle$，即\cref{eq:heavy-nrqcd-decay}\cite[式~(2.13)--(2.14)]{BodwinBraatenLepage1995}。湮灭是局域的：它产生至少两个能量 $\sim m_Q$ 的硬胶子（或轻夸克），湮灭点只在 $\sim1/m_Q$ 的精度内确定\cite[第~II.E~节]{BodwinBraatenLepage1995}。
(2)~系数由在壳 $Q\bar Q$ 散射振幅在 QCD 与 NRQCD 中的匹配逐阶确定，与偶素态无关（例见\cref{prop:heavy-nrqcd-matching}之后的计算）。
(3)~拓扑因子化\cite[第~V.A~节]{BodwinBraatenLepage1995}：对轻强子末态求和时，末态部分子辐射的软、共线发散按 Kinoshita--Lee--Nauenberg（KLN）定理相消；接到重夸克线上的软胶子，若所接传播子离壳 $\order{m_Q}$，红外发散被这一“控制动量”截断，故红外发散部分总可沿离壳 $\ll m_Q$ 的重夸克传播子与短程部分分开，吸收进 NRQCD 矩阵元。例如只含色单态的旧计算中 $\chi_{QJ}$ 衰变的结合能对数\cite{Barbieri:1975am}被同阶的色八重态矩阵元吸收\cite{Bodwin:1992ye}。
(4)~按速度标度规则每阶 $v$ 只有有限个算符（此步只给出 E-S 型陈述）。
\end{proofsketch}

\paragraph{地位与缺口}
在微扰论中为全阶论证（[A]）：拓扑因子化的全阶形式只以概要给出，截断标度的依赖要在 $f_n$ 与 $\langle\mathcal O_n\rangle$ 之间逐阶相消（E-S 型陈述）。BBL 认为其严格程度与“the proofs of factorization in deep-inelastic lepton-hadron scattering”（深度非弹散射中的因子化证明）相当\cite[第~VI.A~节]{BodwinBraatenLepage1995}；就微扰内容而言这与 [A] 一致，在强子层次上则不然：DIS 的矩经色散关系与欧氏区的 OPE 相联系（\cref{ch:dis}），而偶素衰变率是固定物理点（$\sim2m_Q$）上的 Minkowski 量，没有可向欧氏区延拓的变量；BBL 也指出，在共振区对能量抹平的办法“is not available to us in the case of quarkonium annihilation”（对偶素湮灭不可用）\cite[第~II.E~节]{BodwinBraatenLepage1995}。因此其非微扰应用需要局域对偶（\cref{conj:ee-duality}(b)，[C]），不能说“NRQCD 衰变因子化与 OPE 同样严格”。升级途径：证明拓扑因子化的全阶形式；非微扰地需控制局域对偶。

\begin{proposition}[匹配系数 \tagR]\label{prop:heavy-nrqcd-matching}
（$R_{\mathrm{PT}}$：在重整化微扰论中对红外正规化的部分子态严格。）在树图与单圈（$\order{\as}$）阶，电磁衰变的匹配系数如下（本命题只涉及部分子态与微扰的 $Q\bar Q$ 态）。
(a)~（树图）$\Imag f_{ee}({}^3S_1)=\frac{\pi e_Q^2\alpha^2}{3}$，$\Imag f_{\gamma\gamma}({}^1S_0)=\pi e_Q^4\alpha^2$。对相对波函数在原点的值为 $\psi(0)=R(0)/\sqrt{4\pi}$ 的色单态 $S$ 波 $Q\bar Q$ 态，领头阶 $\langle\mathcal O_1\rangle=2N_c|\psi(0)|^2$，于是
\begin{equation}\label{eq:heavy-EMwidths}
\begin{gathered}
\Gamma({}^3S_1\to e^+e^-)=\frac{16\pi N_c}{3}\,\frac{e_Q^2\alpha^2|\psi(0)|^2}{M^2}=\frac{N_ce_Q^2\alpha^2|R(0)|^2}{3m_Q^2},\\
\Gamma({}^1S_0\to\gamma\gamma)=16\pi N_c\,\frac{e_Q^4\alpha^2|\psi(0)|^2}{M^2}=\frac{4N_ce_Q^4\alpha^2|R(0)|^2}{M^2},\qquad M=2m_Q .
\end{gathered}
\end{equation}
% VF-NRQ-2, VF-NRQ-3
(b)~（单圈）阈值处矢量流 $\bar Q\gamma^iQ$ 匹配到 $c_v\,\chi^\dagger\sigma^i\psi$，
\begin{equation}\label{eq:heavy-cv}
c_v=1-2C_F\frac{\as}{\pi}+\order{\as^2},\qquad
\Imag f_{ee}({}^3S_1)=\frac{\pi e_Q^2\alpha^2}{3}\Big(1-4C_F\frac{\as}{\pi}\Big),
\end{equation}
% VF-NRQ-4
故 $\Gamma(V\to e^+e^-)\propto1-\frac{16\as}{3\pi}$（$N_c=3$）。
\uses{def:heavy-nrqcd,def:prelim-msbar}
\end{proposition}

\begin{proof}
(a)~在阈值 $p=\bar p=(m_Q,\bm0)$ 处，入射 $Q(p)\bar Q(\bar p)$ 的自旋投影 $\sum_{s\bar s}\langle\frac12s\,\frac12\bar s|SS_z\rangle\,u(p,s)\bar v(\bar p,\bar s)$ 为
$\Pi_S=\frac{1}{2\sqrt2\,m_Q}(\sla p+m_Q)\,\Gamma_S\,(\sla p-m_Q)$（$\Gamma_0=\gamma_5$，$\Gamma_1=\sla\varepsilon$；$\bar uu=2m_Q$），$\Tr[\bar\Pi_S\Pi_S]=-4m_Q^2$。以相对论归一化的束缚态与色单态波函数 $\delta_{ij}/\sqrt{N_c}$ 组合自由 $Q\bar Q$ 态，领头阶振幅为 $\mathcal M=\frac{\sqrt{2M}}{2m_Q}\sqrt{N_c}\,\psi(0)\Tr[\Gamma_f\Pi_S]$（$\frac{1}{2m_Q}$ 来自两个 $1/\sqrt{2E}$）。对 ${}^3S_1\to\gamma^*\to\ell^+\ell^-$，$\Tr[\gamma^\mu\Pi_1]=-2\sqrt2\,m_Q\varepsilon^\mu$；轻子张量对偶素极化平均为 $\frac13(-g_{\mu\nu}+\frac{P_\mu P_\nu}{M^2})\Tr[\sla k_1\gamma^\mu\sla k_2\gamma^\nu]=\frac43M^2$；乘以 $e^4e_Q^2$、光子传播子的平方 $1/M^4$、两体相空间 $\frac1{8\pi}$ 与 $\frac{1}{2M}$，得 $\Gamma=\frac{N_ce^4e_Q^2|\psi(0)|^2}{3\pi M^2}$，即\cref{eq:heavy-EMwidths}第一式。对 ${}^1S_0\to\gamma\gamma$，两个 Compton 型图之和为 $\Tr[(\gamma^\nu S_Q(p-k_1)\gamma^\mu+\gamma^\mu S_Q(p-k_2)\gamma^\nu)\Pi_0]=\ii\sqrt2\,\epsilon^{\mu\nu\alpha\beta}k_{1\alpha}k_{2\beta}/m_Q^2$（$S_Q$ 为去掉因子 $\ii$ 的夸克传播子），对光子极化求和给出 $4(k_1\cdot k_2)^2/m_Q^4$，计入全同光子的 $\frac12$ 即得第二式。选择定则也直接成立：自旋单态不与单光子耦合，自旋三重态不衰变为两个光子。由 $\Gamma=2\Imag f\,\langle\mathcal O_1\rangle/m_Q^2$ 读出 $\Imag f$，与\cite[式~(A31)、(A34)]{BodwinBraatenLepage1995}一致（BBL 的“$M$”是 $m_Q$，不是偶素质量）；QED 极限 $N_c=e_Q=1$ 给出仲正电子素的 $\alpha^5m/2$。

(b)~严格阈值处 NRQCD 一侧的单圈修正是无标度积分，为零；QCD 一侧取 Feynman 规范，维数正规化同时正规化紫外与红外发散，故 $c_v-1=\delta_V+\delta Z_2^{\rm OS}$（两条外线各贡献 $\frac12\delta Z_2^{\rm OS}$）。顶点修正为
\begin{equation}\label{eq:heavy-vertex}
\Lambda^\mu=-\ii g^2C_F\int\frac{\dd^Dk}{(2\pi)^D}\,\frac{\gamma^\alpha(-\sla p-\sla k+m_Q)\gamma^\mu(\sla p-\sla k+m_Q)\gamma_\alpha}{k^2\,[(p-k)^2-m_Q^2]\,[(p+k)^2-m_Q^2]} ,
\end{equation}
因子 $-\ii g^2C_F$ 来自两个顶点 $(\ii g\gamma t^a)^2$、两个夸克传播子的 $\ii^2$、胶子传播子的 $-\ii$ 与 $t^at^a=C_F$。由 $p^2=m_Q^2$，三个分母为 $k^2$、$k^2\mp2p\cdot k$；Feynman 参数化后令 $k=l+ap$（$a=x_1-x_2$ 为两个夸克分母的参数之差，测度化为 $\int_{-1}^1\dd a\,(1-|a|)$），分母化为 $(l^2-a^2m_Q^2)^3$。用 $\Pi_1$ 投影并作张量约化 $l^\mu l^\nu\to l^2g^{\mu\nu}/D$ 后，分子等于树图的迹乘以 $r_0+r_2l^2$，$r_0=-2m_Q^2[2-(1-\varepsilon)a^2]$，$r_2=\frac{2(1-\varepsilon)^2}{2-\varepsilon}$。做 $D$ 维动量积分，再用 $\int_{-1}^1\dd a\,(1-|a|)|a|^{p}=\frac{2}{(p+1)(p+2)}$ 并按 $\MSbar$ 记 $a_s\equiv\frac{\as C_F}{4\pi}$，得到对 $\varepsilon$ 精确的
\begin{align}
\delta_V&=a_s\,\ee^{\varepsilon\gamma_E}\Big(\frac{\mu^2}{m_Q^2}\Big)^{\varepsilon}\Big\{\Gamma(\varepsilon)\frac{1-\varepsilon}{1-2\varepsilon}+\Gamma(1+\varepsilon)\Big[\frac{2}{\varepsilon(1+2\varepsilon)}-\frac{1}{1-2\varepsilon}\Big]\Big\}\notag\\
&=a_s\Big(\frac3\varepsilon-4+3\ln\frac{\mu^2}{m_Q^2}\Big)+\order{\varepsilon},\label{eq:heavy-dV}\\
\delta Z_2^{\rm OS}&=-a_s\,\ee^{\varepsilon\gamma_E}\Big(\frac{\mu^2}{m_Q^2}\Big)^{\varepsilon}\Gamma(\varepsilon)\,\frac{3-2\varepsilon}{1-2\varepsilon}
=-a_s\Big(\frac3\varepsilon+4+3\ln\frac{\mu^2}{m_Q^2}\Big)+\order{\varepsilon},\label{eq:heavy-dZ2}
\end{align}
% VF-NRQ-4
其中 $\delta_V$ 的第一项来自 $r_2l^2$、第二项来自 $r_0$，$\delta Z_2^{\rm OS}$ 由自能在 $\sla p=m_Q$ 处的导数得到。二者之和精确地为
\begin{equation}\label{eq:heavy-cv-exact}
c_v-1=-\frac{8a_s\,\Gamma(1+\varepsilon)\,\ee^{\varepsilon\gamma_E}}{1-4\varepsilon^2}\Big(\frac{\mu^2}{m_Q^2}\Big)^{\varepsilon}\ \xrightarrow{\ \varepsilon\to0\ }\ -8a_s=-2C_F\frac{\as}{\pi},
\end{equation}
有限且与 $\mu$ 无关。由\cref{eq:heavy-EMwidths}，$\Gamma\propto|c_v|^2=1-4C_F\frac{\as}{\pi}+\order{\as^2}$，与\cite[式~(A37)]{BodwinBraatenLepage1995}一致（BBL 引用\cite{Barbieri:1975ki,Celmaster:1978yz}）。迹、张量约化与参数积分均由符号计算完成并与显式 $4\times4$ 矩阵核对（脚本 \texttt{s5\_\allowbreak cv.py}、\texttt{s7\_\allowbreak extra.py}、\texttt{s8\_\allowbreak review.py}）。
\end{proof}

\begin{example}[阈值处的单圈 $c_v$]\label{ex:heavy-cv}
\cref{eq:heavy-dV,eq:heavy-dZ2}的极点来源不同：顶点的 $\frac3\varepsilon$ 中，$r_2l^2$ 项给出 $\frac1{\varepsilon_{\rm UV}}+1$（取 $\mu=m_Q$），$r_0$ 项给出 $\frac{2}{\varepsilon_{\rm IR}}-5$，后者来自 $a\to0$，即胶子动量趋于零、两条重夸克线同时近壳的区域；$\delta Z_2^{\rm OS}$ 的极点为 $-\frac1{\varepsilon_{\rm UV}}-\frac2{\varepsilon_{\rm IR}}$。紫外极点相消体现 Ward 恒等式 $Z_1=Z_2$；红外极点也相消，且\cref{eq:heavy-cv-exact}与 $\mu$ 无关，说明 NRQCD 矢量流在单圈没有反常维数。稍高于阈值（每个重夸克在质心系中的速度为 $\beta$），完整顶点还含势区域的 Coulomb 奇异性：以胶子质量作红外正规化数值计算单圈阿贝尔顶点（再作 $\alpha\to C_F\as$），在 $\beta=0.01$--$0.08$ 上拟合，振幅修正为 $\frac{C_F\as}{\pi}K(\beta)$，
\begin{equation}\label{eq:heavy-Kbeta}
K(\beta)=\frac{\pi^2}{4\beta}-2+\order{\beta}\qquad(\text{拟合值 }2.46740/\beta-2.0003,\ \pi^2/4=2.46740).
\end{equation}
第一项是 Coulomb 交换，由 NRQCD 的 Coulomb 势复现（平方后即 Sommerfeld 因子 $1+\frac{\pi C_F\as}{2\beta}+\cdots$）；常数 $-2$ 恰为硬区域的 $(c_v-1)/(C_F\as/\pi)$（脚本 \texttt{s6\_\allowbreak threshold.py}）。修正很大：$\as=0.25$ 时 $16\as/(3\pi)\approx0.42$。
\uses{prop:heavy-nrqcd-matching}
\end{example}

%%% BEGIN fig:nrqcd
\begin{figure}[htbp]
\centering
\begin{tikzpicture}[fd,
  t3 heavy/.style={draw,solid,line width=1.3pt,fd arrow=#1},
  t3 heavy/.default=0.5,
  t3 soft/.style={gluon,draw=green!45!black},
  t3 lab/.style={font=\small},
  t3 note/.style={font=\footnotesize}]
  % ===================== (a) colour-singlet channel =====================
  \begin{scope}
    \node[hard blob,circle,minimum size=1.0cm,inner sep=0pt] (H) at (-3.4,0) {};
    \node[soft blob,ellipse,minimum width=2.6cm,minimum height=1.3cm,inner sep=0pt] (L) at (0.8,0) {};
    \draw[gluon] (-5.0,0.85) -- (H.150);  \draw[gluon] (-5.0,-0.85) -- (H.210);
    \draw[gluon] (H.45) -- (-2.15,1.3);
    \draw[t3 heavy] (-2.92,0.17) -- (-0.4,0.17);
    \draw[t3 heavy] (-0.4,-0.17) -- (-2.92,-0.17);
    \draw[bound state] (L.east) -- (3.55,0);
    \node[soft blob,ellipse,minimum width=2.6cm,minimum height=1.3cm,inner sep=0pt] at (0.8,0) {};
    \node[hard blob,circle,minimum size=1.0cm,inner sep=0pt] at (-3.4,0) {};
    \node[t3 lab,left] at (-5.0,0.85) {$g$};  \node[t3 lab,left] at (-5.0,-0.85) {$g$};
    \node[t3 lab,right] at (-2.15,1.3) {$g$};
    \node[t3 lab] at (-1.38,0.5) {$Q\bar Q[{}^3S_1^{[1]}]$};
    \node[t3 note,align=center] at (0.8,0) {$\langle\mathcal O^{H}({}^3S_1^{[1]})\rangle$};
    \node[t3 lab,right] at (3.55,0) {$H\ (J/\psi)$};
    \node[t3 note] at (-3.4,-0.78) {$\sim1/m_Q$};
    \node[t3 note] at (0.8,-0.92) {$\sim1/(m_Qv)$};
    \node[t3 lab,anchor=west] at (-5.6,-1.45) {(a)\ 色单态：短距离过程须多放出一个硬胶子};
  \end{scope}
  % ===================== (b) colour-octet channel =====================
  \begin{scope}[yshift=-3.6cm]
    \node[hard blob,circle,minimum size=1.0cm,inner sep=0pt] (H) at (-3.4,0) {};
    \node[soft blob,ellipse,minimum width=2.6cm,minimum height=1.3cm,inner sep=0pt] (L) at (0.8,0) {};
    \draw[plain] (-5.0,0.85) -- (H.150);  \draw[plain] (-5.0,-0.85) -- (H.210);
    \draw[t3 heavy] (-2.92,0.17) -- (-0.4,0.17);
    \draw[t3 heavy] (-0.4,-0.17) -- (-2.92,-0.17);
    \draw[t3 soft] (1.25,0.55) -- (2.1,1.45);
    \draw[t3 soft] (1.25,-0.55) -- (2.1,-1.45);
    \draw[bound state] (L.east) -- (3.55,0);
    \node[soft blob,ellipse,minimum width=2.6cm,minimum height=1.3cm,inner sep=0pt] at (0.8,0) {};
    \node[hard blob,circle,minimum size=1.0cm,inner sep=0pt] at (-3.4,0) {};
    \node[t3 lab,left] at (-5.0,0.85) {$i$};  \node[t3 lab,left] at (-5.0,-0.85) {$j$};
    \node[t3 lab] at (-1.38,0.5) {$Q\bar Q[n^{[8]}]$};
    \node[t3 note,align=center] at (0.8,0) {$\langle\mathcal O^{H}(n^{[8]})\rangle$};
    \node[t3 lab,right] at (3.55,0) {$H$};
    \node[t3 note,right,text=green!35!black] at (2.1,1.45) {软胶子，$E\sim m_Qv^2$};
    \node[t3 note] at (-3.4,-0.78) {$\sim1/m_Q$};
    \node[t3 lab,anchor=west] at (-5.6,-1.75) {(b)\ 色八重态：颜色在长距离上由软胶子带走};
  \end{scope}
\end{tikzpicture}
\caption{重夸克偶素产生的 NRQCD 因子化图像。灰色团为短距离过程：在距离 $\sim1/m_Q$ 上产生处于确定色、自旋与轨道角动量态 $n={}^{2S+1}L_J^{[1,8]}$ 的 $Q\bar Q$ 对（粗线；箭头为费米子数），
其截面 $\dd\hat\sigma(ij\to Q\bar Q[n]+X)$ 可按 $\alpha_s(m_Q)$ 微扰计算；绿色团为长距离矩阵元 $\langle\mathcal O^{H}(n)\rangle$，描述 $Q\bar Q[n]$（相对速度 $v$，尺度 $\sim1/(m_Qv)$）演化为强子 $H$ 的概率，按 $v$ 的幂次计数（$c\bar c$：$v^2\approx0.3$；$b\bar b$：$v^2\approx0.1$）。
(a)~色单态：例如 $gg\to Q\bar Q[{}^3S_1^{[1]}]+g$，在 $v$ 的领头阶 $Q\bar Q$ 直接成为 $J/\psi$；为满足色与 $C$ 宇称守恒，短距离过程必须多放出一个硬胶子。
(b)~色八重态（如 ${}^1S_0^{[8]}$、${}^3S_1^{[8]}$、${}^3P_J^{[8]}$）：$Q\bar Q$ 对须放出能量 $\sim m_Qv^2$ 的软胶子才能成为色单态强子，相应矩阵元相对于 (a) 有额外的 $v$ 压低。
产生截面的因子化形式 $\dd\sigma(H)=\sum_n\dd\hat\sigma(Q\bar Q[n])\,\langle\mathcal O^H(n)\rangle$ 是猜想（[C]），对衰变的类似公式有全阶论证（[A]）；八重态算符须以 Wilson 线补全规范不变性（见正文）。}
\label{fig:nrqcd}
\end{figure}
%%% END fig:nrqcd
\begin{conjecture}[产生的 NRQCD 因子化 \tagC]\label{conj:heavy-nrqcd-production}
对一切以 $m_Q$ 为硬标度的单举重夸克偶素产生过程，对 $v$ 逐阶，
\begin{equation}\label{eq:heavy-nrqcd-prod}
\sigma(H+X)=\sum_n\frac{F_n(\Lambda)}{m_Q^{d_n-4}}\,\langle0|\mathcal O^H_n(\Lambda)|0\rangle,\qquad
\mathcal O^H_n=\chi^\dagger\mathcal K_n\psi\Big(\sum_X\sum_{m_J}|H+X\rangle\langle H+X|\Big)\psi^\dagger\mathcal K'_n\chi ,
\end{equation}
其中 $F_n$ 是产生 $Q\bar Q[n]$ 的短程截面（可按 $\as(m_Q)$ 计算，与 $H$ 无关），$\mathcal K_n$、$\mathcal K'_n$ 为色与自旋矩阵（高维算符中含协变导数），长程矩阵元普适，且到 $v$ 的任一给定阶只有有限个\cite[第~VI.A~节，式~(6.4)]{BodwinBraatenLepage1995}。
\uses{def:heavy-nrqcd}
\end{conjecture}

\paragraph{地位与缺口}
猜想（[C]），不能说“NRQCD 产生的因子化已证明”。指定末态含 $H$ 使 KLN 求和中缺了一些割：BBL 的论证是，两端都接在重夸克线上的软胶子虽无 KLN 相消，却属于矩阵元；连接与偶素同向的部分子和其他喷注的软胶子，其红外发散“至少在微扰论中”在单举产生率中相消。对强子引起的过程，BBL “merely assume that the Glauber divergences cancel”（仅仅假设 Glauber 发散相消），并认为产生的论证不如衰变严格，与 Drell--Yan 的因子化证明相当\cite[第~VI.A~节]{BodwinBraatenLepage1995}。Nayak、Qiu、Sterman（NQS）在大 $\pT$ 胶子碎裂中发现 NNLO（两圈）有正比于 $\as^2v^2/\varepsilon$ 的不相消红外极点，常规八重态矩阵元不能匹配；规范不变性要求把八重态产生算符的色指标用延伸到无穷远的伴随 Wilson 线 $\Phi_l$ 补全（NQS 为确定起见取 $l$ 为类光矢量，在这一阶结果与 $l$ 的方向无关），补全后的矩阵元恰好产生同样的极点，因子化在 NNLO、$\order{v^2}$ 恢复\cite{NayakQiuSterman2005a,NayakQiuSterman2005b}，对八重态到单态的跃迁推广到 $v$ 的所有阶\cite{NayakQiuSterman2006}（[P]）；若无特殊相消，更高阶可能需要速度展开之外的非局域算符\cite{NayakQiuSterman2005a}。因此不能说“NQS 的规范完备 Wilson 线在 NNLO 不改变结果”：它们在 NNLO 是必需的。两圈协变计算与新进展见第~24~章的\cref{prop:front-nrqcd-ir}（[P]）与\cref{sec:front-quarkonium}。长程矩阵元的普适性与偏振在唯象上存在长期张力\cite{Butenschoen:2010rq,Butenschoen:2011yh,Butenschoen:2012px,Chao:2012iv}（综述\cite{BrambillaEtAl2011,Lansberg:2019adr}），本书不作裁断。升级途径：全阶证明或显式反例。

\begin{proposition}[大 $\pT$ 重夸克偶素产生 \tagA]\label{prop:heavy-quarkonium-fragmentation}
在硬碰撞中产生横动量 $\pT\gg m_Q$ 的重夸克偶素 $H$，其单举截面在 $m_Q^2/\pT^2$ 的领头与次领头幂次因子化为
\begin{equation}\label{eq:heavy-frag}
\dd\sigma_{H}=\sum_f\dd\hat\sigma_{f}\Big(\frac{\pT}{z},\mu\Big)\otimes D_{H/f}(z,\mu)
+\sum_{[Q\bar Q(\kappa)]}\dd\hat\sigma_{[Q\bar Q(\kappa)]}\Big(\frac{\pT}{z},\mu\Big)\otimes\mathcal D_{H/[Q\bar Q(\kappa)]}(z,\mu)
+\order{\frac{m_Q^4}{\pT^4}}\dd\sigma_H .
\end{equation}
第一项（领头幂次）是单部分子碎裂，$D_{H/f}$ 是\cref{def:pdf-ff}意义下的碎裂函数，$\dd\hat\sigma_f$ 是产生部分子 $f$ 的短程截面（初态部分子分布已含在内）；第二项（相对 $m_Q^2/\pT^2$ 压低）是硬碰撞中直接产生的、处于色与自旋态 $\kappa$ 的共线 $Q\bar Q$ 对的碎裂，$\mathcal D$ 为重夸克对（双部分子）碎裂函数，$\otimes$ 还包含对对内两个夸克光锥动量分数的卷积。碎裂函数不作进一步展开，特别是不作 NRQCD 展开。
非微扰地位：\cref{asm:intro-bridge}（MA1、MA2）。
% errtype: E-S
\uses{def:pdf-ff,asm:intro-bridge}
\end{proposition}

\begin{proofsketch}
(1)~领头区域是高能部分子碎裂产生 $H$：标准的共线因子化，截面 $\propto\pT^{-4}$。(2)~在次领头幂次，另一个区域是 $Q\bar Q$ 对在硬碰撞中直接产生且彼此共线；它相对 $m_Q^2/\pT^2$ 压低，但可被 $v$ 的负幂次增强\cite{Fleming:2012wy}，在 NRQCD 中估计其碎裂函数可以重现色单态模型 NLO 计算中的大部分增强\cite{Kang:2011mg}。(3)~对每个区域给出红外安全硬部分的投影算符、碎裂函数的算符定义与演化方程\cite{Kang:2014tta}，在 SCET 中以双部分子碎裂函数写出对所有阶的因子化定理\cite{Fleming:2012wy}（此步只给出 E-S 型陈述）。
\end{proofsketch}

\paragraph{地位与缺口}
全阶论证（[A]），在所引文献（Kang、Qiu、Sterman 及 Ma\cite{Kang:2011mg,Kang:2014tta}；Fleming 等人\cite{Fleming:2012wy}）的层次上。最弱的环节：在强子对撞中，它与一切有色末态的单举截面一样依赖于\cref{conj:dy-coloured-final}（[C]）所涉及的 Glauber 相消；没有初态强子的 $e^+e^-$ 湮灭见\cref{thm:coll-sia}（[A]）。碎裂函数的 NRQCD 因子化与\cref{conj:heavy-nrqcd-production}同为猜想（[C]），不属于本命题。升级途径：对有色末态证明 Glauber 相消。偶素产生与次领头幂次端点因子化的新进展见\cref{sec:front-quarkonium}与\cref{sec:front-nlp}。

% ---------------------------------------------------------------------
\begin{tcolorbox}[enhanced,breakable,title=常见误解,fonttitle=\bfseries,
  colback=gray!3!white,colframe=gray!55!black,boxrule=0.6pt,arc=2pt,
  left=5pt,right=5pt,top=3pt,bottom=3pt,before skip=6pt,after skip=6pt]
\textbf{“$B\to\pi\pi$ 的 QCD 因子化已对所有阶证明。”}不对：领头幂次公式在 $\order{\as}$ 已确立，$T^{\rm I}$ 到两圈显式红外有限，SCET 中有全阶推导，但粲企鹅与硬旁观者项的分歧未决，没有公认的全阶证明（\cref{prop:heavy-qcdf-lp}，[P]）；[A] 只属于 $B\to D\pi$ 与重到轻形状因子。\par
\textbf{“BBNS 公式只在 $\order{\as}$ 成立。”}低估了现状：它是作为对 $\as$ 所有阶的领头幂次陈述提出的，高阶检验与 SCET 推导都支持它。\par
\textbf{“$X_A$、$X_H$ 的参数化是因子化定理的一部分。”}不对：它们参数化发散的卷积（\cref{prop:heavy-qcdf-endpoint}，[V]），是模型。\par
\textbf{“NRQCD 衰变因子化与 OPE 同样严格。”}不对：衰变率是固定物理点上的 Minkowski 量，非微扰地需要局域对偶（\cref{conj:ee-duality}(b)）。\par
\textbf{“NRQCD 产生的因子化已经证明；NQS 的 Wilson 线在 NNLO 不改变结果。”}都不对：产生因子化是猜想（\cref{conj:heavy-nrqcd-production}），规范完备的 Wilson 线在 NNLO 是必需的。
\end{tcolorbox}

\paragraph{本章小结。}
$B\to D\pi$ 与大反冲下的重到轻形状因子在 SCET 中有全阶论证（[A]），前提是模式完备与（对 $B\to D\pi$）端点压低；两个轻介子末态的 BBNS 公式在 $\order{\as}$ 已确立、$T^{\rm I}$ 到两圈显式红外有限，全阶证明尚无共识（[P]）；其幂次修正因卷积发散不能按标准共线形式因子化（[V]），$X_{A,H}$ 只是模型。单举 $B$ 衰变率没有 $\Lqcd/m_b$ 项，端点区是硬函数、喷注函数与形状函数的卷积（[A]）。NRQCD 衰变公式在微扰论中有全阶论证（[A]），电磁匹配系数到单圈严格（[R]）；产生的 NRQCD 因子化是猜想（[C]），在 NNLO 需要规范完备的 Wilson 线；大 $\pT$ 偶素的碎裂因子化有全阶论证，其碎裂函数的 NRQCD 展开不在其中。用于固定运动学点上的强子衰变率，以上结果都还需要局域对偶（[C]）。
